% Mini FSD: CS5330 final report (IEEE conference format). Due 2026-12-08.
% Build: pixi run report (tectonic). Numbers come from the harness log via
% minifsd.eval.report, never typed by hand; figures from docs/report/figures.
% Each section's comment says what it holds and where its numbers come from.
\documentclass[conference]{IEEEtran}
\usepackage[T1]{fontenc}
\usepackage{graphicx}
\usepackage{booktabs}
\usepackage{amsmath}
\usepackage{url}

\begin{document}

\title{Mini FSD: Camera-Only Indoor Navigation\\
Trained on Synthetic Data, Tested on Real Photos}

\author{\IEEEauthorblockN{Oluwasegun Adewola}
\IEEEauthorblockA{CS5330 Pattern Recognition and Computer Vision\\
Northeastern University}}

\maketitle

\begin{abstract}
% 150-200 words, written last: problem, approach (sim data engine + model
% comparison under an edge budget), the headline numbers (E2c floor IoU on
% real rooms, detection mAP, latency), and the main lesson.
\end{abstract}

\section{Introduction}
% The robot, the camera-only constraint, why synthetic data (labels are free
% and exact), why an edge budget (8 GB M1, 2 threads, 50 ms for perception).
% Contributions as a short list: the data engine, the shared harness, the E1
% to E3 comparisons, the furniture finding (D11, E2c).

\section{Related Work}
% Short: domain randomization (Tobin et al.), foundation features for
% transfer (DINOv2/DINOv3), real-time detectors (YOLO26, RF-DETR), efficient
% segmentation (SegFormer), monocular depth (Depth Anything V2 and 3).

\section{System}
% MuJoCo world from generated floor plans (R1 to R3), ROS 2 and Nav2 on top
% (R5 to R9), ground-truth pose (ADR 0004). One figure: the pipeline.

\section{Synthetic Data Engine}
% D1 to D12: cameras (height, pitch, field of view), looks (floors, walls,
% glass, lighting), signs, doors, camera effects, poses that see open floor,
% boxes for detection, COCO people, and furniture with floor-like textures.
% Split by floor plan. v9 and v10 sizes and box counts from the plan's log.

\section{Evaluation Harness and Edge Budget}
% One harness for every candidate: same splits, metrics (floor IoU, mAP at
% IoU 0.5, AbsRel and delta1 with per-image and global fits) and latency on
% the M1 under a 2-thread cap. Pareto front; the pick rule. Latency
% measurement conditions (one cool session) and the untimed DA3 case.

\section{Experiments}

\subsection{E2: Floor Segmentation}
% SegFormer-B0 fine-tuned against frozen DINOv2 ViT-S/14 with a linear probe,
% trained on v9, then both on v10 (E2c). Table from minifsd.eval.report:
% nyuv2-val, v9 val, v10 val IoU, precision and recall, p50 latency. The
% figure: accuracy against latency. Overlays of the counter/bed failure.

\subsection{E1: Detection}
% YOLO26n against RF-DETR Nano at 320 px on v10 plus COCO people; mAP at 0.5
% per class on v10 val and COCO people val; latency through ONNX Runtime.

\subsection{E3: Depth}
% Depth Anything V2 Small (V5) against DA3 Small (accuracy only, run on
% Colab) on NYUv2 test and v9 val; E3b if it lands.

\subsection{Real Photos}
% Floor 15 for model choice, floors 14 and 2 used once for final numbers
% (P1 to P5). Fallback stated if campus access slipped.

\section{Discussion}
% What transferred and what didn't: the sim-val versus real gap, furniture,
% glass, latency over budget and what E6 must do. Limitations.

\section{Conclusion}
% Three or four sentences.

\bibliographystyle{IEEEtran}
\bibliography{refs}

\end{document}
